% Introduction to Part XI. Unnumbered text: no \label except the \phantomsection target.
\phantomsection\label{intro:XI}
The flow of Part~X bounds its coefficients, $|\beta_j|\le C_\beta\mathsf K$, with $\mathsf K$ the bound on the generations (Remark~\ref{rem:flow-constants}), but gives no sign and no lower bound for them (Theorem~\ref{thm:N}(iii)), and it assumes no asymptotic behaviour of the couplings (Remark~\ref{rem:sf-flow}(a)). The ultraviolet behaviour is the statement that $g_k^{-2}-g_K^{-2}$ grows like $(K-k)\beta_\infty$, with $\beta_\infty>0$. Part~XI computes the one-loop part of the coefficient, summed over the steps, and transfers it to the flow. As in Part~X, the statements about couplings concern the small-field flow on a fixed torus; nothing is claimed about large fields or a continuum limit.

\paragraph{Why the one-loop coefficient suffices.} The coefficient $\beta_{j+1}$ of the flow has two parts (Section~\ref{sec:ol-uv}).
\begin{itemize}
\item Its one-loop part is the coefficient of the one-loop term $\Gamma^{\rm step}_j=-\frac12\log\det\Delta_j+\sum_c\log\mathfrak j_c$ of the step from scale $j$, a function of the background alone (Section~\ref{sec:ol-composite}, paragraph ``The one-loop terms of the steps''), up to an error exponentially small in the side of the torus at scale $j+1$. This is (F2) of Proposition~\ref{prop:flow-F}; it rests on the torus independence of the one-loop pieces (Proposition~\ref{prop:one-loop-local}, Lemma~\ref{lem:OL}).
\item The rest of $\beta_{j+1}$ is of order $g_j^2p_0(g_j)^6$, by parity in the coupling; this is (F3).
\end{itemize}
So the couplings follow the one-loop recursion up to errors that add up to a power of $\log(K-k+2)$.

\paragraph{What this Part proves.}
\begin{itemize}
\item \emph{One function} (Sections~\ref{sec:ol-composite} and~\ref{sec:ol-chart}). The one-loop term $\Gamma^{(1)}_k$ of the composite $k$-step map, given by Laplace's method (Proposition~\ref{prop:ol-laplace}), is the sum of the $\Gamma^{\rm step}_j$ up to constants (Proposition~\ref{prop:ol-composition}). By the symmetry of Ba\l aban's averaging, the marginal tensors of these functions are multiples of $\delta_{\mu\nu}|p|^2-p_\mu p_\nu$, with no term linear in $p$ (Lemma~\ref{lem:ol-symmetry}). In four dimensions the change of scale preserves marginal tensors, so the tensors of the steps add up to that of $\Gamma^{(1)}_k$ (Corollary~\ref{cor:ol-cumulative}). In the Landau chart, $\Gamma^{(1)}_k$ is given by the determinants of $\bL_*$, of $\bQ_k\bG_*\bQ_k^T$ and of a ghost operator (Theorem~\ref{thm:ol-chart}).
\item \emph{The computation} (Sections~\ref{sec:ol-vacpol} and~\ref{sec:ol-feynman}). Uniformly in $k$, the marginal tensor of the vacuum polarization at $U=1$ equals that of the lattice Feynman-gauge operators up to $O(1)$: the averaging adds operators of finite rank per cell, or raises the power-counting degree (Proposition~\ref{prop:ol-reduction}, Lemmas~\ref{lem:ol-power-counting} and~\ref{lem:ol-remainders}). The Feynman-gauge coefficient is $k\beta_\infty+O(1)$, with $\beta_\infty=\frac{11C_A}{24\pi^2T_\rho}\log L$ (Theorem~\ref{thm:ol-log}). Hence $\sum_{j<k}T^{(1)}_{j+1}=(k\beta_\infty+O(1))(\delta_{\mu\nu}|p|^2-p_\mu p_\nu)$ (Theorem~\ref{thm:ol-cumulative}).
\item \emph{The ultraviolet behaviour} (Section~\ref{sec:ol-uv}). With (F1)--(F3), for every small $g$ and every $K$ there is a bare coupling for which the flow stays in $(0,\gamma]$, has $g_K=g$, and satisfies $|g_k^{-2}-g^{-2}-(K-k)\beta_\infty|\le C(\log(K-k+2))^{6q+1}$, where $p_0(g)=(\log g^{-1})^q$ (Theorem~\ref{thm:ol-uv}). This is the analogue of Ba\l aban's (0.31) for the couplings of this flow, except that for the last $m_1$ steps its lower half holds up to an additive constant (Remark~\ref{rem:ol-remains}).
\end{itemize}
These results form Theorem~\ref{thm:M}.

\paragraph{Off the main line.} Section~\ref{sec:gen-locality} proves that, at fixed couplings and at flat configurations, the generations converge as the torus grows (Theorem~\ref{thm:gen-locality}). With it, the $\beta_j$ converge to the coefficient that Ba\l aban's definition \cite[(1.20)--(1.22)]{BalYM4I} assigns to the generations of this note (Proposition~\ref{prop:beta-identification}). No rate is proved, and Theorem~\ref{thm:M} does not use these results.

\paragraph{Imports.} Parts~I--VII (in particular the operators at $U=1$ of Part~I), Propositions~\ref{prop:one-loop-polymer} and~\ref{prop:one-loop-local} and Lemma~\ref{lem:local-inventory} of Part~IX, and from Part~X Propositions~\ref{prop:jacobians} and~\ref{prop:first-step} and Lemma~\ref{lem:first-step-formula}; Sections~\ref{sec:ol-uv} and~\ref{sec:gen-locality} also use the flow, Section~\ref{subsec:torus-independence}, Theorem~\ref{thm:weak-bg} and the graded domains of Part~VIII (Definition~\ref{def:graded}); Section~\ref{sec:gen-locality} uses in addition the weakened families of Part~IX (Lemmas~\ref{lem:weak-basic} and~\ref{lem:weak-chain}, Theorem~\ref{thm:weak-inverse}, Corollary~\ref{cor:weak-covariance}).

\paragraph{Details for this Part.} The companion has two chapters on Part~XI. The first expands the composite one-loop term, the marginal tensors, the infinite-volume operators, the Landau chart and the vacuum polarization; the second the lattice Feynman-gauge coefficient, the period density, the properties (F1)--(F3), the ultraviolet limit and the generation locality.
